\section{Background and Related Work}
\label{sec:background}
\label{sec:related}

This section first reviews public tomato maturity datasets and the detectors
trained on them. It then explains detector evaluation and the distinction
between a count and a detection score. The remaining subsections review work on
the factors varied in this study: annotation conventions, robustness to
lighting, occlusion, and inference on an edge device. Table~\ref{tab:prior}
summarises the resulting research gap.

\subsection{Tomato maturity datasets and detectors}
\label{sec:datasets-bg}

Public annotated data for tomato maturity is recent. The two datasets used in
this work were both released in 2020. \emph{Laboro Tomato} was published by
Laboro.AI, a Japanese company, and contains 804 images collected at a local
farm with two different cameras, at 3024$\times$4032 and 3120$\times$4160
pixels \cite{laborotomato2020}. It has bounding boxes and instance masks for
six classes, which combine two fruit sizes with three maturity stages. The
stages are defined by the share of the fruit surface that is red: at least
90\% for fully ripened, 30--89\% for half-ripened, and less than 30\% for green,
which follows the logic of the USDA colour classes \cite{usda1991tomato}.
\emph{tomatOD} was published by Tsironis \emph{et al.}\ \cite{tsironis2020tomatod}
and contains 277 images taken over three days in April 2019 in a soilless
greenhouse in Crete, Greece. The authors used a 12-megapixel machine-vision
camera to simulate a camera mounted on a robot arm. Two agriculturists labelled
its 2,418 boxes as unripe, semi-ripe or fully ripe. They deliberately left fruit
unlabelled when more than 50\% was occluded or when it was beyond the second
plant row. Thus, the datasets were collected in different countries with
different cameras, annotators, and definitions of ripeness.

Other public tomato datasets exist but are less suited to maturity counting.
The Rob2Pheno dataset of Afonso \emph{et al.}\ \cite{afonso2020tomato} contains
123 RGB-D images (720$\times$1280) from a commercial greenhouse in the
Netherlands. Its red and green labels were derived automatically from a colour
ratio and were not annotated by hand. The TomatoPlantFactory dataset
\cite{wu2023plantfactory} contains 520 images of the Micro-Tom variety under
artificial lighting in China, with two classes (green and red). A frequently
used Kaggle dataset \cite{makeml2020tomato} has 895 images with a single
``tomato'' class and no maturity labels. Laboro Tomato and tomatOD are, to our
knowledge, the only two public datasets with three hand-annotated maturity
stages in commercial-style greenhouses, which is why this work uses them.

Researchers have applied several detector architectures to these datasets since
2020. YOLO-Tomato (2020) adapted YOLOv3 with dense blocks and circular boxes
to detect tomatoes under changing light and occlusion, without classifying maturity
\cite{liu2020yolotomato}. Mu \emph{et al.}\ (2020) detected highly occluded
green tomatoes with Faster R-CNN and estimated their number
\cite{mu2020occluded}, and Afonso \emph{et al.}\ (2020) detected and counted
tomatoes with Mask R-CNN \cite{afonso2020tomato}. Later work focused on
maturity: Su \emph{et al.}\ (2022) classified four maturity classes with a
MobileNet-based YOLOv3 \cite{su2022seyolov3}; Li \emph{et al.}\ (2023) improved
YOLOv5 for four maturity stages \cite{li2023yolov5tomato}; and a second Li
\emph{et al.}\ (2023) added attention to YOLOv8 for maturity detection and
counting \cite{li2023mhsayolov8}. Camacho and Morocho-Cayamcela (2023) compared
Mask R-CNN and YOLOv8 on a three-class version of Laboro Tomato
\cite{camacho2023tomato}. We replicated their protocol as the starting point of
this work. Each study trains and tests on a single dataset, usually its own, and
reports mAP or precision and recall at one setting. We found no study that
trains on Laboro Tomato and tests on tomatOD, or the reverse. The effect of the
different annotation conventions has therefore not been measured before.

\subsection{Maturity is a colour judgement, and its boundaries are conventions}

A tomato ripens through a continuous change in pigment, from green chlorophyll
through orange carotenoids to red lycopene. The change is gradual, so the
discrete labels green, half-ripened and fully ripened divide a continuum. The
fruit has no physical marker at the boundary between stages.
Postharvest science measures this change with colour indices, most often the
CIELAB $a^*$ value or the $a^*/b^*$ ratio \cite{lopezcamelo2004color}, and
standards for grading tomatoes by colour define stages by the share of the
surface that has turned red \cite{usda1991tomato}.

A detector trained to name a stage learns the boundaries used by the people who
annotated its training data. When it is evaluated against labels drawn by
another team, its error therefore includes the disagreement between the two
annotation conventions. Because the boundary is a colour threshold,
illumination, camera white balance and over-exposed highlights also move a
fruit across it. Section~\ref{sec:colour} measures where the two public datasets
place these boundaries.

\subsection{How detectors are evaluated}
\label{sec:metrics-bg}

Most fruit detectors use a single-stage design: the network predicts a set of
boxes, each with a class and a confidence score, in one forward pass. This
design is generally less accurate than two-stage methods such as Mask R-CNN,
but is fast enough for real-time use \cite{redmon2016yolo}. The YOLO family is
therefore widely used in agricultural vision. This work uses YOLO11
\cite{khanam2024yolov11, jocher2024ultralytics}, whose five sizes (n, s, m, l,
x) share one architecture and differ only in network width and depth. They
therefore form a controlled series for measuring the effect of extra capacity.
Before the boxes are returned, non-maximum suppression (NMS) removes duplicate
boxes for the same object, keeping the most confident one \cite{neubeck2006nms}.

Most studies compare detectors using mean average precision (mAP), as defined
for the COCO benchmark \cite{lin2014coco}. A predicted box counts as
correct when its intersection over union (IoU) with an annotated box exceeds a
threshold. IoU is the ratio between the overlapping area of the two boxes and
the total area covered by either box. Precision and recall are computed as the
confidence threshold is lowered through the ranked list of predictions. The
area under the precision--recall curve is averaged over classes and over IoU thresholds from
0.5 to 0.95 (mAP50--95), or taken at IoU 0.5 only (mAP50). Two properties of
mAP matter for this paper. First, it averages over all confidence thresholds and
describes how well the predictions are ranked. It does not describe the model's
behaviour at the single threshold used in an application. Second, it is computed
per box and therefore does not directly measure a count derived from the boxes.

\subsection{Counting is a different task from detecting}

A greenhouse application needs a fruit count. To produce one, the application
counts every prediction above a chosen confidence threshold. The number of
predictions above this threshold depends on the calibration of the model's
confidence scores. Calibration describes how closely a score of 0.7
corresponds to a 70\% chance of being correct, and modern networks are known to be poorly
calibrated \cite{guo2017calibration}. Detection errors can also cancel in the
count. An image in which one fruit is missed and one false fruit is detected has
a perfect count but two detection errors.

Detection quality and counting quality have been shown to diverge in practice.
In a comparison of fruit detection and counting methods for yield mapping in
apple orchards, the method with the best detection result differed from the
method with the best count \cite{hani2020comparative}. Reporting mAP without
the count therefore leaves the error that matters to the application
unmeasured.

\subsection{Robustness to dataset and lighting changes}

Computer-vision models often perform worse on a dataset other than the one used
for training. Datasets can differ enough for a classifier to identify the source
of an image, and models trained on one dataset can transfer poorly to others
\cite{torralba2011unbiased}. Robustness to common
image changes, such as noise, blur, weather and changes in brightness, has
since been benchmarked systematically for image classification, and models
that rank well on clean images can rank differently under such changes
\cite{hendrycks2019benchmarking}. White-balance errors in particular have been
shown to reduce the accuracy of trained networks \cite{afifi2019whatelse}.
Data augmentation addresses this problem by adding synthetic variation that the
training set lacks \cite{shorten2019augmentation}. In the YOLO training
pipeline, the colour part of this augmentation is a random change of hue,
saturation and brightness within three fixed limits \cite{bochkovskiy2020yolov4,
ultralytics2026cfg}; Section~\ref{sec:hsv} tests whether widening those limits helps.

Agricultural images are especially exposed to lighting change. Greenhouse light
varies with the time of day, the weather and the glazing, and the feature being
classified, colour, is the one that lighting changes most directly. Previous
agricultural work has addressed illumination mainly by changing the input
representation or by collecting more varied data. Kirk \emph{et al.}\ replace
the RGB input with colour-opponent features and report better generalisation to
a real strawberry farm than an RGB detector of comparable accuracy
\cite{kirk2020lab}. Haggag \emph{et al.}\ vary dataset size, camera viewpoint
and illumination for a tomato greenhouse detector and find illumination the
hardest of the three shifts to generalise across
\cite{haggag2024generalization}. Neither study tests the colour augmentation
already present in the training pipeline, and neither reports whether the
differences it observes remain when the model is retrained with a different
random seed. None of the tomato maturity studies reviewed in
Section~\ref{sec:datasets-bg} evaluates its detector under controlled changes
of exposure or white balance.

\subsection{Evaluation under occlusion}

Occlusion has been studied through part-based architectures, controlled
benchmarks and synthetic augmentation. Context-aware CompositionalNets show
that the context around an object becomes especially important when the object
is strongly occluded \cite{wang2020robust}. Gilroy \emph{et al.}\ argue that
broad categories such as ``partly occluded'' hide useful information, and
evaluate seven detectors at measured occlusion levels from 0\% to 99\%
\cite{gilroy2023benchmark}. Kassaw \emph{et al.}\ compare real and artificial
occlusion and find that performance depends on the type of occluder
\cite{kassaw2026occlusion}. During training, methods such as Cutout, Random
Erasing and CutMix hide or replace rectangular image regions to make networks
less dependent on any single part of an object \cite{devries2017cutout, zhong2020randomerasing,
yun2019cutmix}.

In agriculture, Liao \emph{et al.}\ place real leaves over apples and control
the share of each fruit that is hidden \cite{liao2025amodal}. This work follows
the same principle, a known hidden fraction and an occluder taken from the same
scene, but uses a rectangular strip of background, and adds two controls: a
flat black strip, to measure the effect of the occluder's appearance, and the
natural overlap between fruit, which needs no synthetic occluder at all.

\subsection{Running the model on an edge device}

A harvesting or scouting aid is used in the field, so its speed must be
measured on its intended edge device. Running a model in the cloud requires
a reliable network connection and adds a network round trip to every frame,
which is often slower and less predictable than local inference
\cite{kang2017neurosurgeon, shi2016edge}. Reviews of edge intelligence in agriculture identify
connectivity as the main constraint on practical systems and argue for
on-device inference \cite{shaheen2023edge}. Deep learning on smartphones is
recent: Core ML and the first iPhone Neural Engine appeared in 2017
\cite{apple2017ios11, apple2017iphonex}, and benchmarks of neural network inference on phones date
from 2018--2019 \cite{ignatov2019aibenchmark}.

Moving a model from a training workstation to a smartphone changes several
parts of the inference pipeline at once. The weights are usually converted to
16-bit floating point; NMS and image
preprocessing may be built into the exported model; execution moves to a
vendor-specific neural accelerator with its own numerical behaviour; and long
runs heat the device, which may cause the operating system to reduce its speed.
In Core ML, the term compute units denotes the processors that the
runtime is allowed to use, such as the CPU and Neural Engine. A latency figure
can therefore only be interpreted together with the device, the permitted
compute units, the operating-system version, and whether the run was short or
sustained.

Edge deployment in agriculture has mostly been studied on single-board
computers and dedicated accelerators. Mazzia \emph{et al.}\ run a reduced
YOLOv3 for apple detection on a Raspberry Pi with a neural compute stick and on
two NVIDIA Jetson boards, and treat the frame rate on these boards as the main
constraint \cite{mazzia2020edge}. Work on citrus compares architectures by
their speed on the target hardware \cite{dasilva2023citrus}. Model scaling has
been studied as a design problem for YOLO detectors, varying width, depth and
input resolution together against a hardware budget \cite{wang2021scaled}.
For tomatoes specifically, Zeng \emph{et al.}\ (2023) quantised a
lightweight YOLOv5 and ran it in an Android application at 26.5 frames per
second \cite{zeng2023mobile}; TSE-YOLO (2025) runs ripeness detection and
segmentation on an Android phone \cite{jia2025tseyolo}; GAE-YOLO (2025) runs on
an NVIDIA Jetson TX2 with a maturity assessment \cite{liu2025gaeyolo}; and Kafi
\emph{et al.}\ (2025) run a YOLOv8 tomato detector, without maturity stages, on
the Neural Engine of an iPhone receiving video from a drone
\cite{kafi2025drone}.
These studies report a throughput figure for one platform. This work keeps one
exported model fixed and measures its latency across devices, compute-unit
settings, operating-system versions and run lengths.

\subsection{What a single benchmark number hides}

Reported improvements vary across training runs. Retraining the same
configuration with a different random seed changes the data order, the
augmentation and the initial weights, and the resulting spread is often as
large as the differences that papers claim; taking that variance into account
changes which comparisons hold up \cite{bouthillier2021variance}. Results from a
single training run are still the norm in applied agricultural vision, which is
why this work repeats the colour-augmentation experiments over three seeds and
uses paired statistical tests (Section~\ref{sec:stats}).

\subsection{Summary and research gap}

Table~\ref{tab:prior} summarises which of the factors studied here previous
tomato work has addressed. Each column is a factor that can change a reported
counting error without any change to the model. Most previous studies report
accuracy on one dataset with one model; counting, deployment on a device and
seed variance each appear in only a few, and no study varies the annotation
convention, the confidence threshold, or the lighting in a controlled way. This
work evaluates a standard detector and measures how much each factor changes
the result. The contribution is this controlled evaluation; no new detector is
proposed.

\begin{table*}[t]
\centering
\caption{Factors addressed by previous tomato detection studies and by this
work. \checkmark: measured or reported; ($\checkmark$): partly, e.g.\ discussed
qualitatively or for one model only; --: not addressed. ``Edge'' means that
latency was measured on the target device.}
\label{tab:prior}
\setlength{\tabcolsep}{4pt}
\begin{tabular}{lcccccccc}
\toprule
Study & Year & Maturity & Count & Cross-dataset & Occlusion & Lighting & Threshold & Edge \\
\midrule
Liu \emph{et al.}\ \cite{liu2020yolotomato}        & 2020 & --         & --         & -- & ($\checkmark$) & ($\checkmark$) & -- & -- \\
Mu \emph{et al.}\ \cite{mu2020occluded}            & 2020 & --         & \checkmark & -- & \checkmark     & --             & -- & -- \\
Afonso \emph{et al.}\ \cite{afonso2020tomato}      & 2020 & ($\checkmark$) & ($\checkmark$) & -- & --     & --             & -- & -- \\
Su \emph{et al.}\ \cite{su2022seyolov3}            & 2022 & \checkmark & --         & -- & --             & --             & -- & -- \\
Li \emph{et al.}\ \cite{li2023yolov5tomato}        & 2023 & \checkmark & --         & -- & --             & --             & -- & -- \\
Li \emph{et al.}\ \cite{li2023mhsayolov8}          & 2023 & \checkmark & \checkmark & -- & --             & --             & -- & -- \\
Camacho \emph{et al.}\ \cite{camacho2023tomato}    & 2023 & \checkmark & --         & -- & --             & --             & -- & -- \\
Zeng \emph{et al.}\ \cite{zeng2023mobile}          & 2023 & --         & --         & -- & --             & --             & -- & \checkmark \\
Haggag \emph{et al.}\ \cite{haggag2024generalization} & 2024 & --      & --         & ($\checkmark$) & --     & \checkmark     & -- & -- \\
Jia \emph{et al.}\ \cite{jia2025tseyolo}           & 2025 & \checkmark & --         & -- & --             & --             & -- & \checkmark \\
Kafi \emph{et al.}\ \cite{kafi2025drone}           & 2025 & --         & --         & -- & --             & --             & -- & \checkmark \\
\midrule
This work                                          & 2026 & \checkmark & \checkmark & \checkmark & \checkmark & \checkmark & \checkmark & \checkmark \\
\bottomrule
\end{tabular}
\end{table*}
